% GENERATED FILE. Edit canonical lessons, then run npm run generate:books.
% canonical-source-hash: fnv1a64:331ba9501bb6666f
% canonical-lessons: IT-C268-andata-e-ritorno, IT-C268-ho-perso, IT-C268-ripasso-stazione

\chapter{A Return Ticket --- And When It Goes Wrong}
\label{ch:it-a-return-ticket}
\hlchaptermodality{268}

\begin{chapteropening}
\textbf{By the end of this chapter:} \emph{I can buy a ticket, one way or return, and report a problem such as a missed train or connection.}

Complete IT-C268-ripasso-stazione: one morning at the station, from finding it to missing the train and asking what next.
\end{chapteropening}

\section[andata e ritorno]{andata e ritorno — the going and the coming back}

\label{lesson:IT-C268-andata-e-ritorno}



\begin{quote}
Say \textquotedblleft{}a ticket\textquotedblright{}. (\emph{Un biglietto.}) Ask for a coffee the polite way. (\emph{Vorrei un caffè.})

Put the two together and you are standing at a ticket window.
\end{quote}

\subsection*{You'll want to know: andata e ritorno}
\textbf{andata e ritorno} — \textquotedblleft{}there and back\textquotedblright{}: a return ticket. \textbf{solo andata} — \textquotedblleft{}one way\textquotedblright{}, the going only.

\begin{quote}
\emph{Un biglietto di andata e ritorno.} — A return ticket. \emph{Solo andata.} — One way only.
\end{quote}

\textbf{ritorno} you have, \textquotedblleft{}the return\textquotedblright{}. \textbf{andata} is new, and you can almost hear how it was made: it is \emph{andato}, \textquotedblleft{}gone\textquotedblright{}, turned into a feminine noun, \textquotedblleft{}the going\textquotedblright{}. \textbf{solo} means \textquotedblleft{}only\textquotedblright{}, the same as \emph{soltanto}.

\begin{grammarlens}[title={Grammar Lens: vorrei un biglietto per…}]
The whole request takes one frame:

\begin{quote}
\textbf{Vorrei un biglietto per} Roma, \textbf{andata e ritorno}. — I'd like a return ticket to Rome.
\end{quote}

\textbf{per} plus the place names where you are going, the same \emph{per} as in \emph{per andare}. No article before a city: \emph{per Roma}, \emph{per Milano}.

Then the questions that come back at you, and the answers you need:

\begin{quote}
— \emph{Andata e ritorno?} — \emph{Sì, andata e ritorno.} / \emph{No, solo andata.} — \emph{Che binario?} — \emph{Binario tre.}
\end{quote}

\emph{Che binario?}, \textquotedblleft{}which platform?\textquotedblright{}, is the same \emph{che} that asks \emph{che cosa?}.
\end{grammarlens}

\begin{cousinweb}
\textbf{andata} is built on \textbf{andare}, whose own origin is not settled: EtymDB 2.1 gives it no Latin parent, and the candidates usually proposed, such as Latin \emph{ambulāre}, \textquotedblleft{}to walk\textquotedblright{}, are disputed. \textbf{ritorno} is \emph{ri-}, \textquotedblleft{}back\textquotedblright{}, plus \emph{tornare}, \textquotedblleft{}to turn\textquotedblright{}, inherited from Latin \emph{tornāre}, \textquotedblleft{}to turn on a lathe\textquotedblright{} (Lewis \& Short, under \emph{torno}); EtymDB 2.1 derives the noun through French \emph{retour}. Either way, a return is a turning round.
\end{cousinweb}

\subsection*{Your turn}
Say these aloud:

\begin{itemize}
\raggedright
  \item \emph{andata e ritorno}, then \emph{solo andata}
  \item \emph{Vorrei un biglietto per Roma, andata e ritorno.}
  \item the one-way version — \emph{Vorrei un biglietto per Milano, solo andata.}
  \item ask the platform — \emph{Scusi, che binario?}
\end{itemize}

\begin{tcolorbox}[breakable,colback=teal!4,colframe=teal!35!black,title={Before you move on}]
What is a one-way ticket? (\emph{\textbf{Solo andata.}}) Which word goes before the city, and does it take an article? (\emph{\textbf{Per}} — and \textbf{no} article.) What is \emph{andata}, taken apart? (\emph{\textbf{Andato}}, \textquotedblleft{}gone\textquotedblright{}, made a noun: \textbf{the going}.)
\end{tcolorbox}
\section[ho perso il treno]{ho perso il treno — when it goes wrong}

\label{lesson:IT-C268-ho-perso}



\begin{quote}
Ask for a one-way ticket to Rome. (\emph{Vorrei un biglietto per Roma, solo andata.})

Now the ticket is in your pocket, and the train has gone without you.
\end{quote}

\subsection*{You'll want to know: ho perso il treno}
\textbf{ho perso il treno} — \textquotedblleft{}I've missed the train\textquotedblright{}.

You have \textbf{perdere}, \textquotedblleft{}to lose\textquotedblright{}. Italian uses the same verb for a train you missed: you \textbf{lost} it. Its past does not follow the \emph{-uto} pattern of \emph{credere}, \emph{creduto}. The participle is short: \textbf{perso}.

\begin{quote}
\textbf{Ho perso il treno.} — I've missed the train. \textbf{Ho perso la coincidenza.} — I've missed my connection. \textbf{Ho perso il biglietto.} — I've lost my ticket.
\end{quote}

Italian says \emph{il} biglietto where English says \emph{my} ticket: whose it is goes without saying.

\begin{grammarlens}[title={Grammar Lens: reporting a problem}]
A complaint to a stranger has three parts, and the first is a courtesy:

\begin{quote}
\textbf{Scusi, ho un problema.} — Excuse me, I have a problem. \textbf{Ho perso il treno.} — What went wrong. \textbf{C'è un altro treno oggi? A che ora?} — What you need.
\end{quote}

Or the fault is not yours: \emph{Il treno è in ritardo}, \textquotedblleft{}the train is late\textquotedblright{}. Say what happened, in the passato prossimo, then ask the question that gets you moving again.
\end{grammarlens}

\begin{cousinweb}
\textbf{perdere} is inherited from Latin \emph{perdere} (EtymDB 2.1), built on \emph{per-}, \textquotedblleft{}through\textquotedblright{}, and \emph{dare}, \textquotedblleft{}to give\textquotedblright{}: to give something away for good, to squander, to lose (Lewis \& Short, under \emph{perdo}). Its Latin participle was \emph{perditus} (Lewis \& Short), but Italian made new ones of its own: the short \textbf{perso}, and the longer \emph{perduto}, which sounds older and more literary. Neither source says how \emph{perso} was formed, so that step is left open here.
\end{cousinweb}

\subsection*{Your turn}
Say these aloud:

\begin{itemize}
\raggedright
  \item \emph{Ho perso il treno.} — then \emph{Ho perso la coincidenza.}
  \item the three parts — \emph{Scusi, ho un problema. Ho perso il treno. C'è un altro treno oggi?}
  \item the fault that is not yours — \emph{Il treno è in ritardo.}
  \item the worst one — \emph{Ho perso il biglietto.}
\end{itemize}

\begin{tcolorbox}[breakable,colback=teal!4,colframe=teal!35!black,title={Before you move on}]
What is the past participle of \emph{perdere}? (\emph{\textbf{Perso.}}) How do you say you missed your connection? (\emph{\textbf{Ho perso la coincidenza.}}) What are the three parts of a complaint? (A \textbf{courtesy}, \textbf{what went wrong}, \textbf{what you need}.) And what are the two Latin pieces of \emph{perdere}? (\emph{\textbf{Per-}}, through, and \emph{\textbf{dare}}, to give.)
\end{tcolorbox}
\section[(review)]{Review — one morning at the station}

\label{lesson:IT-C268-ripasso-stazione}



\begin{quote}
Say \textquotedblleft{}I've missed the train\textquotedblright{}. (\emph{Ho perso il treno.})
\end{quote}

\begin{grammarlens}[title={What you've built}]
Two chapters of travel, and they make one morning. You find the station, you buy the ticket, something goes wrong, and you ask your way out of it.

Every step is a sentence you have said: \emph{Scusi, per andare alla stazione?} · \emph{Vorrei un biglietto per Roma, andata e ritorno.} · \emph{Scusi, ho un problema: ho perso il treno.} · \emph{C'è un altro treno oggi?}

And when you tell someone about it afterwards, the morning becomes a story: \emph{Il treno era in ritardo, quindi ho perso la coincidenza.} Scene, arrow, event.
\end{grammarlens}

\subsection*{Your turn}
Say these aloud:

\begin{itemize}
\raggedright
  \item find it — \emph{Scusi, per andare alla stazione?}
  \item the answer, said back — \emph{Dritto, poi a sinistra?}
  \item buy it — \emph{Vorrei un biglietto per Roma, andata e ritorno.}
  \item the problem — \emph{Scusi, ho un problema. Ho perso il treno.}
  \item tell it afterwards — \emph{Il treno era in ritardo, quindi ho perso la coincidenza.}
\end{itemize}

\begin{tcolorbox}[breakable,colback=teal!4,colframe=teal!35!black,title={Before you move on}]
What is a return ticket? (\emph{\textbf{Andata e ritorno.}}) What do you say before the problem itself? (\emph{\textbf{Scusi, ho un problema.}}) In \emph{il treno era in ritardo, quindi ho perso la coincidenza}, which verb is the scene? (\emph{\textbf{Era}}.)
\end{tcolorbox}
